% Part V -- Geospatial Data and Applied Mapping
% Chapter 14: GeoPandas foundations

\h{GeoPandas Foundations}
\label{ch:14}

Many design and planning questions include a location: Which buildings lie
inside a service area? How far is a transit stop from a studio? Which features
intersect a study boundary? GeoPandas extends a pandas DataFrame with spatial
geometry and a coordinate reference system (CRS), while retaining familiar
table operations and Matplotlib plotting.

\begin{NoteBox}
\textbf{Prerequisites.} You should be comfortable with DataFrames, column
selection, functions, file paths, and the Figure--Axes pattern from
\booklink{ch:13}{Chapter 13}. Spatial analysis also requires careful attention
to units and coordinate systems.
\end{NoteBox}

\begin{tcolorbox}[title=Learning outcomes,colframe=orange,colback=white,arc=2mm]
After completing this chapter, you should be able to:
\begin{itemize}
    \item explain how a GeoDataFrame extends a DataFrame;
    \item create Point geometry from longitude and latitude columns;
    \item inspect geometry types, bounds, and CRS metadata;
    \item distinguish declaring a CRS with \c{set_crs()} from transforming
          coordinates with \c{to_crs()};
    \item create distance-based buffers in a projected CRS;
    \item layer spatial data on one Matplotlib Axes; and
    \item export a reproducible site-service map from a complete script.
\end{itemize}
\end{tcolorbox}

\hh{The GeoDataFrame mental model}

A regular DataFrame stores observations in rows and variables in columns. A
GeoDataFrame adds one active \c{geometry} column and CRS metadata. Each geometry
may be a Point, LineString, or Polygon. Together, geometry and CRS tell Python
both \emph{where} a feature is and how its coordinate numbers relate to Earth.

\bookfigure[0.94\textwidth]{Figures/Generated/ch14_geodataframe_model.pdf}
{A GeoDataFrame combines attribute columns, one active geometry column, and coordinate-reference-system metadata.}
{fig:ch14-geodataframe-model}

The attribute columns still support pandas operations such as filtering,
sorting, grouping, and aggregation. GeoPandas adds spatial operations such as
buffering, testing intersections, and spatially joining two layers.

\begin{SyntaxBox}{Three components work together}
\begin{itemize}
    \item \textbf{Attributes} describe each feature, such as its name or capacity.
    \item \textbf{Geometry} stores the feature's spatial form and coordinates.
    \item \textbf{CRS} defines the coordinate system and therefore the units and
          geographic interpretation of those coordinates.
\end{itemize}
A geometry without a correct CRS can be plotted, but its distances, alignment,
and location may be wrong.
\end{SyntaxBox}

\hh{Install once and import when needed}

Install the mapping libraries in the same Python environment used to run your
scripts.

\begin{verbatim}
python -m pip install geopandas geodatasets matplotlib
\end{verbatim}

The conventional imports keep the longer library names readable.

\begin{lstlisting}[
    language=python,
    style=block,
    caption={Import the table, spatial, and plotting libraries},
    label={c14_imports},
    captionpos=b
]
import geopandas as gpd
import matplotlib.pyplot as plt
import pandas as pd
\end{lstlisting}

The alias \c{gpd} identifies GeoPandas, \c{pd} identifies pandas, and \c{plt}
identifies Matplotlib's plotting interface. An import makes a library available;
it does not load a dataset or create a map.

\hh{Create points from longitude and latitude}

Suppose a table contains one row per site and two numeric coordinate columns.
Create the ordinary DataFrame first, then derive Point geometry.

\begin{lstlisting}[
    language=python,
    style=block,
    caption={Convert a site table into a GeoDataFrame},
    label={c14_create_points},
    captionpos=b
]
import geopandas as gpd
import pandas as pd

sites = pd.DataFrame({
    "name": ["Main Hall", "Design Studio", "Transit Stop"],
    "longitude": [-76.1355, -76.1318, -76.1372],
    "latitude": [43.0386, 43.0371, 43.0410],
})

points = gpd.GeoDataFrame(
    sites,
    geometry=gpd.points_from_xy(sites["longitude"], sites["latitude"]),
    crs="EPSG:4326",
)

print(points[["name", "geometry"]])
print(points.crs)
\end{lstlisting}

\begin{SyntaxBox}{Read the constructor from the inside out}
\begin{enumerate}
    \item \c{sites["longitude"]} supplies horizontal coordinates and
          \c{sites["latitude"]} supplies vertical coordinates. The order is
          \c{x, y}, not latitude first.
    \item \c{gpd.points_from_xy(...)} creates one Point for every row.
    \item \c{geometry=...} chooses that result as the active geometry column.
    \item \c{crs="EPSG:4326"} declares that the source numbers are geographic
          longitude and latitude coordinates.
    \item The final assignment preserves attributes and geometry in \c{points}.
\end{enumerate}
The nested form is compact, but each argument answers a separate question:
which table, which geometry, and which CRS?
\end{SyntaxBox}

\begin{TryBox}{Verify the coordinate order}
Print \c{points.total_bounds}. The four values are minimum x, minimum y,
maximum x, and maximum y. Longitude values near $-76$ and latitude values near
$43$ are plausible for Syracuse. Reversed coordinates would be an immediate
warning sign.
\end{TryBox}

\hh{Inspect before analyzing}

Spatial data deserves the same inspect-first discipline as any table. Check
its rows, schema, geometry, CRS, and extent before calculating or plotting.

\begin{lstlisting}[
    language=python,
    style=block,
    caption={Inspect a spatial table systematically},
    label={c14_inspect},
    captionpos=b
]
print(points.head())
print(points.columns.tolist())
print(points.geometry.geom_type.value_counts())
print(points.crs)
print(points.total_bounds)
\end{lstlisting}

\c{points.geometry} accesses the active geometry GeoSeries.
\c{geom_type} reports the type of each feature, and \c{total_bounds} summarizes
the entire layer as \c{[min_x, min_y, max_x, max_y]}. These checks often reveal
bad coordinate order, a missing CRS, or an unexpected mixture of geometry types.

\hh{CRS: declare correctly, then transform deliberately}

EPSG codes provide compact identifiers for coordinate reference systems.
EPSG:4326 represents longitude and latitude in angular degrees. It is useful
for storage and exchange, but one degree is not a fixed ground distance.
Distance-based work therefore needs an appropriate projected CRS with linear
units. Syracuse lies in UTM zone 18N, for which EPSG:32618 uses meters.

\begin{lstlisting}[
    language=python,
    style=block,
    caption={Transform coordinates before measuring distance},
    label={c14_transform_crs},
    captionpos=b
]
points_utm = points.to_crs("EPSG:32618")

print(points_utm.crs)
print(points_utm.geometry.x.round(1))
print(points_utm.geometry.y.round(1))

service_areas = points_utm.buffer(400)
print(service_areas.area.round(1))
\end{lstlisting}

\begin{SyntaxBox}{\c{set_crs()} and \c{to_crs()} do different jobs}
\begin{itemize}
    \item \c{set_crs(...)} attaches or corrects CRS metadata when the existing
          coordinate numbers are already expressed in that CRS. It does not
          recalculate coordinates.
    \item \c{to_crs(...)} transforms geometry into a different CRS and returns
          new coordinate values.
    \item \c{buffer(400)} uses the current CRS units. After transformation to
          EPSG:32618, the radius is 400 meters.
\end{itemize}
Never use \c{set_crs()} merely to make two layers display the same CRS label.
If both original CRS definitions are correct, transform one layer with
\c{to_crs()} instead.
\end{SyntaxBox}

For more CRS detail, consult the
\href{https://docs.geopandas.org/en/stable/docs/user_guide/projections.html}{official GeoPandas projections guide}.

\hh{Layer spatial data on one Axes}

GeoPandas \c{plot()} is built on Matplotlib. Pass the same \c{ax} to every layer
so that polygons and points share one coordinate space.

\begin{lstlisting}[
    language=python,
    style=block,
    caption={Draw service polygons below site points},
    label={c14_layer_plot},
    captionpos=b
]
fig, ax = plt.subplots(figsize=(7, 5), layout="constrained")

service_areas.plot(
    ax=ax,
    color="tab:blue",
    alpha=0.16,
    edgecolor="tab:blue",
)
points_utm.plot(
    ax=ax,
    color="tab:orange",
    edgecolor="black",
    markersize=45,
    zorder=3,
)

ax.set_aspect("equal")
plt.show()
plt.close(fig)
\end{lstlisting}

The first layer establishes the map extent. The second call reuses \c{ax}; it
does not create a separate figure. Transparency keeps overlapping polygons
visible, and \c{zorder=3} places the points above them. Equal aspect prevents
one map unit from appearing wider than the same unit vertically.

\hh{Applied example: map four 400-meter service areas}
\label{sec:c14-application}

The following program is deliberately complete: it prepares data, constructs
geometry, transforms coordinates, calculates buffers, draws layers, annotates
sites, and exports the finished figure. Reading those stages is as important
as memorizing individual methods.

A runnable copy is provided in \c{examples/c14_site_service_areas.py}.

\begin{lstlisting}[
    language=python,
    style=block,
    caption={Build and export a site-service map},
    label={c14_site_service_application},
    captionpos=b
]
from pathlib import Path

import geopandas as gpd
import matplotlib.pyplot as plt
import pandas as pd

# 1. Store attributes and source coordinates.
sites = pd.DataFrame({
    "name": ["Main Hall", "Design Studio", "Transit Stop", "Library"],
    "longitude": [-76.1355, -76.1318, -76.1372, -76.1286],
    "latitude": [43.0386, 43.0371, 43.0410, 43.0395],
})

# 2. Create point geometry in the source CRS.
points = gpd.GeoDataFrame(
    sites,
    geometry=gpd.points_from_xy(sites["longitude"], sites["latitude"]),
    crs="EPSG:4326",
)

# 3. Transform to a meter-based CRS before measuring.
points = points.to_crs("EPSG:32618")

# 4. Build a polygon GeoDataFrame from 400 m buffers.
service_areas = gpd.GeoDataFrame(
    points[["name"]].copy(),
    geometry=points.buffer(400),
    crs=points.crs,
)

# 5. Draw polygons first and points second.
fig, ax = plt.subplots(figsize=(7.4, 5.2), layout="constrained")
service_areas.plot(
    ax=ax,
    color="tab:blue",
    alpha=0.16,
    edgecolor="tab:blue",
    linewidth=1.2,
)
points.plot(
    ax=ax,
    color="tab:orange",
    edgecolor="black",
    markersize=48,
    zorder=3,
)

# 6. Label each point using its projected x and y coordinates.
for row in points.itertuples():
    ax.annotate(
        row.name,
        (row.geometry.x, row.geometry.y),
        xytext=(5, 5),
        textcoords="offset points",
        fontsize=8,
    )

# 7. Add map context and export.
ax.set(
    title="Four 400 m site-service areas",
    xlabel="Easting (m), UTM zone 18N",
    ylabel="Northing (m), UTM zone 18N",
)
ax.set_aspect("equal")
ax.grid(alpha=0.18)

output_folder = Path("study_figures")
output_folder.mkdir(exist_ok=True)
fig.savefig(
    output_folder / "site_service_areas.png",
    dpi=200,
    bbox_inches="tight",
)
plt.show()
plt.close(fig)
\end{lstlisting}

\bookfigure[0.86\textwidth]{Figures/Generated/c14_site_buffers.pdf}
{Rendered output from Listing~\ref{c14_site_service_application}. The projected axes make the 400-meter radius meaningful.}
{fig:c14-site-buffers}

\begin{SyntaxBox}{Trace the application as a data pipeline}
\begin{enumerate}
    \item The DataFrame preserves the original longitude and latitude columns.
    \item \c{points_from_xy()} creates geometry; the GeoDataFrame constructor
          declares the source CRS.
    \item \c{to_crs()} produces meter-based Point coordinates.
    \item \c{points.buffer(400)} returns a GeoSeries of Polygon geometry.
    \item A second GeoDataFrame associates each service Polygon with a site name.
    \item Repeated \c{plot(ax=ax, ...)} calls compose layers on one Axes.
    \item \c{itertuples()} supplies named attributes and geometry to the label loop.
    \item \c{Path.mkdir()} and \c{fig.savefig()} make the result reproducible.
\end{enumerate}
Each intermediate variable has one spatial meaning. That separation makes
units and geometry easier to inspect when a map looks wrong.
\end{SyntaxBox}

The circles overlap most strongly near the center, suggesting shared coverage.
That observation is descriptive, not yet an accessibility claim: real travel
depends on streets, barriers, entrances, and the meaning assigned to each site.

\hh{Ask spatial questions with predicates}

A spatial predicate returns \c{True} or \c{False} for a geometric relationship.
For example, a point can be within a polygon, and two polygons can intersect.

\begin{lstlisting}[
    language=python,
    style=block,
    caption={Test simple geometric relationships},
    label={c14_predicates},
    captionpos=b
]
first_area = service_areas.geometry.iloc[0]

inside_first_area = points.geometry.within(first_area)
overlaps_first_area = service_areas.geometry.intersects(first_area)

print(points.loc[inside_first_area, "name"])
print(service_areas.loc[overlaps_first_area, "name"])
\end{lstlisting}

\c{within(first_area)} compares every Point in the GeoSeries with one Polygon.
The resulting Boolean Series can filter rows with \c{loc}. Likewise,
\c{intersects(first_area)} identifies service polygons sharing any space with
the first polygon. Predicate direction matters: a point is within a polygon,
whereas the polygon contains the point.

\hh{Common mistakes and useful diagnoses}

\begin{itemize}
    \item \textbf{Latitude and longitude are reversed.} Check \c{total_bounds}
          and remember that Point construction uses x then y.
    \item \textbf{A distance is calculated in degrees.} Transform to a suitable
          projected CRS before buffering or measuring.
    \item \textbf{A CRS label is overwritten.} Use \c{set_crs()} only when
          declaring what the existing numbers already mean.
    \item \textbf{Layers do not align.} Print both CRS values and transform one
          layer to the other layer's CRS.
    \item \textbf{Each layer appears in a different figure.} Create one Axes and
          pass \c{ax=ax} to every GeoPandas plot call.
    \item \textbf{The map is visually distorted.} Use \c{ax.set_aspect("equal")}
          for projected coordinates.
\end{itemize}

\hh{Practice ladder}

\hhh{Level 0 -- Read}

In Listing~\ref{c14_create_points}, identify the attribute columns, geometry
column, and source CRS.

\hhh{Level 1 -- Modify}

Add a fifth site to the application. Re-run the program and confirm that its
Point and service Polygon both appear.

\hhh{Level 2 -- Explain}

Change the buffer radius from 400 to 250 meters. Explain why making this change
before \c{to_crs()} would not represent 250 meters.

\hhh{Level 3 -- Build}

Add a \c{category} column and use two marker shapes or a legend so categories
remain distinguishable without relying only on color.

\hhh{Level 4 -- Challenge}

Create a function \c{make_service_areas(points, radius_m)} that validates a
positive radius, returns a new polygon GeoDataFrame, and leaves the input
GeoDataFrame unchanged.

\hh{Practice solutions}

\textbf{Level 0:} The attributes are \c{name}, \c{longitude}, and \c{latitude};
\c{geometry} is the active geometry column; the source CRS is EPSG:4326.

\textbf{Level 2:} A buffer operates in the current CRS units. Before projection,
EPSG:4326 coordinates are angular degrees, not meters.

\textbf{Level 4}

\begin{lstlisting}[
    language=python,
    style=block,
    caption={One solution for a reusable service-area function},
    label={c14_solution_level4},
    captionpos=b
]
import geopandas as gpd

def make_service_areas(points, radius_m):
    if radius_m <= 0:
        raise ValueError("radius_m must be positive")
    if points.crs is None:
        raise ValueError("points must have a CRS")
    if points.crs.is_geographic:
        raise ValueError("project points to a meter-based CRS first")

    result = points.drop(columns="geometry").copy()
    return gpd.GeoDataFrame(
        result,
        geometry=points.geometry.buffer(radius_m),
        crs=points.crs,
    )
\end{lstlisting}

The guards fail early when the radius, CRS, or units would make the result
misleading. Copying non-geometry attributes keeps the input unchanged, and the
returned GeoDataFrame receives new Polygon geometry.

\hh{Chapter checkpoint}

\begin{enumerate}
    \item What three components give a GeoDataFrame its spatial meaning?
    \item In what order does \c{points_from_xy()} receive longitude and latitude?
    \item What does \c{total_bounds} report?
    \item Why should a 400-meter buffer not be calculated in EPSG:4326?
    \item How does \c{set_crs()} differ from \c{to_crs()}?
    \item Why pass the same \c{ax} to multiple \c{plot()} calls?
    \item What type of result does a spatial predicate produce?
\end{enumerate}

\hh{Checkpoint answers}

\begin{enumerate}
    \item Attribute columns, active geometry, and correct CRS metadata.
    \item Longitude as x, then latitude as y.
    \item Minimum x, minimum y, maximum x, and maximum y for the layer.
    \item EPSG:4326 uses degrees, so 400 would not mean 400 meters.
    \item \c{set_crs()} declares existing coordinates; \c{to_crs()} transforms them.
    \item The layers then share one coordinate space and one map.
    \item A Boolean value per compared geometry, usually stored as a Series.
\end{enumerate}

\begin{tcolorbox}[title=Chapter summary,colframe=orange,colback=white,arc=2mm]
A GeoDataFrame combines familiar table attributes with geometry and a CRS.
Inspect all three, transform to suitable units before measuring, and build maps
as deliberate Matplotlib layers. These foundations prepare you to load
documented datasets and connect layers through spatial joins in
\booklink{ch:15}{Chapter 15}.
\end{tcolorbox}
